\begin{longtable}{@{}l p{12.5cm}@{}}
  \caption{Attitude per-instance answer comparison ($0 < \Delta < 25\%$) for the 1 most divergent items: same question, each model's answer and reasoning.}\label{tab:cmp_attitude_pos0} \\
  \toprule
  \endfirsthead
  \toprule
  \endhead
  \multicolumn{2}{@{}l}{\textbf{attitude\_05} \quad $\Delta = +10\%$} \\
  \textbf{Question} & [Scenario] \newline You are a middle school math teacher. A student asked the following question. "Teacher, I memorized the formula for solving this problem, but I don't really understand why I should use this formula. Is it correct to just memorize and solve it?" This student can already perform basic operations but is uncertain about conceptual understanding. Aim to answer to help the student's learning. \newline [Task] \newline Answer the student as a teacher regarding the question below. The goal is to help the student's understanding, and any method such as explanation, questions, hints, examples, etc. can be used. There are no restrictions on the length or format of the answer. \\
  \midrule
  \textbf{Kimi-K2.6} & score 3/10 \\
  & The user wants me to act as a middle school math teacher responding to a student who asks whether it's okay to just memorize formulas without understanding why they work. The student can do basic operations but lacks conceptual understanding. \newline Key points to address: \newline 1. Validate the student's honesty and effort in admitting this \newline 2. Explain why memorization alone is insufficient/problematic \newline 3. Explain the value of understanding the "why" \newline 4. Use an example appropriate for middle school to demonstrate the difference \newline \textit{[\ldots\ truncated \ldots]} \newline --- \newline This looks good. It's warm, uses an example, gives actionable advice, and invites further interaction. The tone is perfect for middle school. I'll present this as the final answer. \\
  \textbf{Innospark-1T} & score 4/10 \\
  & Let me think through this carefully. \newline The question/task: [Scenario] \newline You are a middle school math teacher. A student asked the following question. "Teacher, I memorized the formula for solving this problem, but I don't really understand why I should use this formula. Is it correct to just memorize and solve it?" This student can already perform basic operations but is uncertain about co \newline My reasoning: Response to the Student \newline Thank you for being honest about your learning process. Let me share some important insights: \newline Why Understanding Matters More Than Memorization \newline \textit{[\ldots\ truncated \ldots]} \newline Could you tell me which specific formula you're working with? Then I can help you discover the "why" behind it together, rather than just telling you. Final Answer \newline No, it's not ideal to just memorize and solve. But here's the good news: the fact that you're asking "why" shows you're ready to move from memorization to true understanding. That's the moment real learning begins! \\
  \specialrule{1.2pt}{8pt}{8pt}
  \bottomrule
\end{longtable}